\section{Experimental Setup}
\label{sec:method}

\subsection{Synthetic Regression Tasks}
\label{sec:tasks}
We generate $N=5{,}000$ evenly spaced inputs $x\in[0,10]$ and three targets of increasing structural
complexity\ifcompact\else\ (\cref{fig:tasks})\fi, each with additive Gaussian noise $\varepsilon\sim\mathcal N(0,\sigma^2)$:
\begin{align}
y_{\text{lin}}   &= 2x + \varepsilon, & \sigma&=1,\\
y_{\text{curv}}  &= 0.5x^2 + \sin x + \varepsilon, & \sigma&=2,\\
y_{\text{chain}} &= s_a + \sin(s_a) + 0.5\,\mathrm{sgn}(\sin x) + \varepsilon, & \sigma&=0.5,
\end{align}
where $s_a = x\cos(0.5x)$. The \emph{chained} target is a composition of three stages: an oscillating
ramp, a nonlinearity applied to that ramp, and a discontinuous square wave. This was designed so that a model whose
structure mirrors the stages could, in principle, utilize it. Because the noise is known, each task has an
irreducible test MSE of $\sigma^2$ (1, 4, and 0.25, respectively), which we use as an absolute reference.
Inputs are scaled to $[0,1]$. Targets are standardized with the training set's mean and standard deviation, the
models are trained on the standardized values, and predictions are mapped back to original units, so every
reported MSE is directly comparable to the noise floors. A single random 80/20
train/test split (seed 42) is shared by every model and task, giving 4,000 training and 1,000 test points.

\ifcompact\else  % dropped in the 6-page build
\begin{figure*}[t]
\centering
\includegraphics[width=\textwidth]{fig_tasks}
\caption{The three synthetic regression tasks ($N=5{,}000$). (a)~linear with $\sigma=1$; (b)~quadratic plus
sinusoid with $\sigma=2$; (c)~a three-stage composition containing a square-wave discontinuity, $\sigma=0.5$.}
\label{fig:tasks}
\end{figure*}
\fi

% Architecture schematics for the six parameter-matched models (Section III).
% FC n = fully connected layer with n outputs; R = ReLU. Edit freely in Overleaf.
\begin{figure*}[t]
\centering
\begin{tikzpicture}[
  font=\scriptsize,
  >={Stealth[length=1.6mm]},
  io/.style={circle, draw=black!70, minimum size=4.2mm, inner sep=0pt},
  fc/.style={rectangle, rounded corners=1pt, draw=black!70, fill=#1, minimum height=4.2mm,
             minimum width=11mm, inner sep=1.5pt, align=center},
  fc/.default=blue!8,
  op/.style={circle, draw=black!70, fill=white, minimum size=4mm, inner sep=0pt},
  mod/.style={draw=black!55, dashed, rounded corners=2pt, inner sep=2.5pt},
  lab/.style={font=\footnotesize\bfseries, anchor=west},
  note/.style={font=\tiny, text=black!60},
  every edge/.style={draw=black!70, ->},
]
%------------------------------------------------ (a) Sequential
\begin{scope}[shift={(0,0)}]
  \node[lab] at (-0.3,1.05) {(a) Sequential \normalfont\scriptsize(385 p)};
  \node[io] (x) at (0,0) {$x$};
  \node[fc] (h) at (1.25,0) {FC 128\\+R};
  \node[fc] (o) at (2.65,0) {FC 1};
  \node[io] (y) at (3.75,0) {$\hat y$};
  \draw (x) edge (h) (h) edge (o) (o) edge (y);
  \node[note] at (1.9,-0.62) {depth 2 $\cdot$ width 128};
\end{scope}
%------------------------------------------------ (b) Parallel
\begin{scope}[shift={(5.9,0)}]
  \node[lab] at (-0.3,1.05) {(b) Parallel \normalfont\scriptsize(385 p)};
  \node[io] (x) at (0,0) {$x$};
  \node[fc=orange!10] (a) at (1.25,0.33) {FC 64+R};
  \node[fc=orange!10] (b) at (1.25,-0.33) {FC 64+R};
  \node[op] (c) at (2.35,0) {\tiny$\Vert$};
  \node[fc=orange!10] (o) at (3.2,0) {FC 1};
  \node[io] (y) at (4.15,0) {$\hat y$};
  \draw (x) edge (a) (x) edge (b) (a) edge (c) (b) edge (c) (c) edge (o) (o) edge (y);
  \node[note] at (2.0,-0.78) {two branches, concatenated};
\end{scope}
%------------------------------------------------ (c) Modular
\begin{scope}[shift={(11.9,0)}]
  \node[lab] at (-0.3,1.05) {(c) Modular \normalfont\scriptsize(385 p)};
  \node[io] (x) at (0,0) {$x$};
  \node[fc=teal!10] (h) at (1.25,0) {FC 128\\+R};
  \node[fc=teal!10] (o) at (2.5,0) {FC 1};
  \node[io] (y) at (3.6,0) {$\hat y$};
  \node[mod, fit=(h)(o), label={[note]below:self-contained module}] {};
  \draw (x) edge (h) (h) edge (o) (o) edge (y);
\end{scope}
%------------------------------------------------ (d) H-A
\begin{scope}[shift={(0,-2.5)}]
  \node[lab] at (-0.3,1.05) {(d) H-A: Parallel of Sequentials \normalfont\scriptsize(385 p)};
  \node[io] (x) at (0,0) {$x$};
  \node[fc=yellow!14] (a1) at (1.1,0.36) {FC 12+R};
  \node[fc=yellow!14] (a2) at (2.45,0.36) {FC 12+R};
  \node[fc=yellow!14] (b1) at (1.1,-0.36) {FC 12+R};
  \node[fc=yellow!14] (b2) at (2.45,-0.36) {FC 12+R};
  \node[op] (c) at (3.45,0) {\tiny$\Vert$};
  \node[fc=yellow!14, minimum width=8mm] (o) at (4.2,0) {FC 1};
  \node[io] (y) at (5.0,0) {$\hat y$};
  \draw (x) edge (a1) (a1) edge (a2) (x) edge (b1) (b1) edge (b2)
        (a2) edge (c) (b2) edge (c) (c) edge (o) (o) edge (y);
  \node[note] at (2.2,-0.8) {two 2-layer branches (depth + breadth)};
\end{scope}
%------------------------------------------------ (e) H-B
\begin{scope}[shift={(6.0,-2.5)}]
  \node[lab] at (-0.3,1.05) {(e) H-B: Staged Modular Chain \normalfont\scriptsize(385 p)};
  \node[io] (x) at (0,0) {$x$};
  \node[fc=magenta!8, minimum width=8mm] (A) at (0.95,0) {$M_A$};
  \node[fc=magenta!8, minimum width=8mm] (B) at (2.05,0) {$M_B$};
  \node[fc=magenta!8, minimum width=8mm] (C) at (3.15,0) {$M_C$};
  \node[op] (s) at (4.2,0) {\tiny$+$};
  \node[io] (y) at (4.9,0) {$\hat y$};
  \draw (x) edge (A) (A) edge (B) (B) edge (C) (C) edge (s) (s) edge (y);
  % outputs a, b also feed later stages and the sum (top); x skips to B and C (bottom)
  \draw[->, black!50] (A.north) to[out=35, in=145, looseness=0.9] (C.north);
  \draw[->, black!50] (A.north east) to[out=30, in=150, looseness=0.8] (s.north);
  \draw[->, black!50] (B.north east) to[out=35, in=150, looseness=0.9] (s.north west);
  \draw[->, black!50] (x.south) to[out=-35, in=-145, looseness=0.9] (B.south);
  \draw[->, black!50] (x.south east) to[out=-30, in=-150, looseness=0.8] (C.south);
  \node[note] at (2.4,-1.05) {$M_{A,B,C}$: FC 33/32/31+R$\rightarrow$FC 1;\; $\hat y = a+b+c$};
\end{scope}
%------------------------------------------------ (f) H-C
\begin{scope}[shift={(12.15,-2.5)}]
  \node[lab] at (-0.3,1.05) {(f) H-C: Trunk + Heads \normalfont\scriptsize(387 p)};
  \node[io] (x) at (0,0) {$x$};
  \node[fc=green!10] (t1) at (1.05,0) {FC 20+R};
  \node[fc=green!10] (t2) at (2.4,0) {FC 15+R};
  \node[fc=green!10, minimum width=7mm] (h1) at (3.5,0.36) {FC 1};
  \node[fc=green!10, minimum width=7mm] (h2) at (3.5,-0.36) {FC 1};
  \node[op] (m) at (4.3,0) {\tiny$\mu$};
  \node[io] (y) at (4.95,0) {$\hat y$};
  \node[mod, fit=(t1)(t2), label={[note]below:shared trunk}] {};
  \draw (x) edge (t1) (t1) edge (t2) (t2) edge (h1) (t2) edge (h2) (h1) edge (m) (h2) edge (m) (m) edge (y);
\end{scope}
\end{tikzpicture}
\caption{The six parameter-matched architectures (385 trainable parameters; 387 for \HC). Top row: the three
single-paradigm baselines from the pilot study. Bottom row: hybrids combining depth, breadth, and modularity.
FC~$n$ = fully connected layer with $n$ outputs, R = ReLU, $\Vert$ = concatenation, $\mu$ = mean.
Note that (a) and (c) compute the same function class; they differ only in how the code is organized,
which lets us use them as a control pair for implementation overhead (Section~\ref{sec:pilot}).}
\label{fig:arch}
\end{figure*}


\subsection{Architectures}
\label{sec:archs}
All six models (\cref{fig:arch}) are one-input, one-output multilayer models with ReLU
activations~\cite{nair2010rectified}, each with exactly 385 trainable parameters, so that differences in accuracy
or speed cannot be attributed to capacity. The one exception is \HC: no width pair for two-head wiring totals
385, so it uses the nearest configuration, 387 (+0.5\%). \Cref{tab:arch} lists their parameter counts,
multiply--accumulate operations (MACs) per sample, critical-path depth (linear layers between input and
output), tensor operations dispatched per forward call, and the low-level (ATen) kernels those operations execute.

\textbf{Single-paradigm baselines.} \emph{Sequential} is a 1--128--1 network. \emph{Parallel} feeds the input
to two 64-unit branches whose activations are concatenated before a linear read-out, in the spirit of
multi-branch designs such as Inception~\cite{szegedy2015going}. \emph{Modular} wraps a hidden layer and read-out
in a self-contained module. With 128 hidden units it calculates exactly the same function class as
Sequential, and under the same seed it initializes to the same weights. We keep it as a
\emph{control}: any measured difference between Sequential and Modular is implementation overhead or
timing noise, never architecture.

\textbf{Hybrids.} \textbf{\HA} (\emph{parallel of sequentials}) runs two 2-layer, 12-unit stacks side by side,
combining depth and breadth. \textbf{\HB} (\emph{staged modular chain}) has three sub-networks $M_A$, $M_B$, $M_C$ (hidden widths 33, 32,
and 31, chosen to total exactly 385 parameters) that mirror the three stages of the chained target: $M_B$ receives $(x,a)$, $M_C$ receives
$(x,a,b)$, and the output is $a+b+c$. This gives densely connected skip links in the spirit of
DenseNet~\cite{huang2017densenet} and an additive output similar to residual learning~\cite{he2016resnet}.
\textbf{\HC} (\emph{trunk + parallel heads}) shares a 2-layer trunk (20 and 15 units) between two linear heads whose
outputs are averaged, the hard-parameter-sharing pattern from multi-task learning~\cite{caruana1997multitask}.
No set of architectures can match parameters, MACs, and operation count at once. Fixing parameters forces the
other two to vary, and that difference is what lets us test which cost measure actually predicts latency.

\begin{table}[t]
\centering
\caption{Architecture cost profile. Depth: linear layers on the longest path. Ops / Kern.: PyTorch operations /
ATen kernels per forward call. Latency: median single-thread CPU time over 10 shuffled repeats; comp.: with
\texttt{torch.compile}.}
\label{tab:arch}
\setlength{\tabcolsep}{2.4pt}
\begin{tabular}{@{}lrrrrrrrr@{}}
\toprule
 & & \textbf{MACs/} & & & & \multicolumn{3}{c}{\textbf{Latency (\textmu s)}}\\
\cmidrule(l){7-9}
\textbf{Model} & \textbf{Params} & \textbf{sample} & \textbf{Depth} & \textbf{Ops} & \textbf{Kern.} & \textbf{b=1} & \textbf{comp.} & \textbf{4096}\\
\midrule
\tablebody{tables/tab_arch_body.tex}
\bottomrule
\end{tabular}
\end{table}

\subsection{Training Protocol}
Every model is implemented and trained in PyTorch~\cite{paszke2019pytorch} with one shared training loop. We use
Adam~\cite{kingma2015adam} (learning rate $10^{-3}$, $\beta=(0.9,0.999)$, $\epsilon=10^{-7}$, matching the Keras
default used in the pilot), mean-squared-error loss, mini-batches of 32 with per-epoch shuffling, and
100 epochs (12,500 optimizer steps). To test whether the chained-task results reflect undertraining, we also
train Sequential, Parallel, \HA, and \HB on the chained task for 300 epochs under the same protocol. Each (model, task) pair is trained with five seeds (0--4) that control
initialization and batch order. We report mean\,$\pm$\,standard deviation across seeds, following
recommendations to account for training variance in benchmark comparisons~\cite{bouthillier2021accounting}.

\subsection{Metrics}
\begin{itemize}
\item \textbf{Test MSE} and $R^2$ on the 1,000 held-out points after the final epoch.
\item \textbf{Convergence epoch}: the first epoch whose test MSE is within 5\% of that run's best.
\item \textbf{Training cost}: wall-clock training time divided by optimizer steps (ms/step).
\item \textbf{Inference latency}: median of 50 forward passes over the full test set (batch 1{,}000).
\end{itemize}
Significance against the Sequential baseline uses Welch's two-sample $t$-test on test MSE ($n=5$ per arm,
$p<0.05$).\ifcompact
 One timing group with a contention outlier was re-measured; accuracy was identical.
\else
 One timing group (\HB, chained task) was re-measured after a single seed's inference time deviated by
30\% from the same network on the other tasks; its accuracy was bit-identical across both passes.
\fi

\subsection{Latency Sweep}
\label{sec:sweep}
Real deployments rarely run inference at exactly the batch size used for evaluation, so we also time every
model at batch sizes $\{1, 8, 64, 512, 1000, 4096\}$ with 1 and 4 CPU threads. Each configuration
gets 50 warm-up calls followed by 1,000 timed calls, and the whole grid is repeated 10 times in a freshly shuffled
order so that run order and cache state cannot result in bias towards any model. We report the median of the per-repeat medians,
with a 95\% interval across repeats.
MACs are counted with forward hooks on every linear layer. Dispatched operations are counted by intercepting
PyTorch's \texttt{\_\_torch\_function\_\_} protocol during one forward pass, and the kernels those operations
execute are counted with \texttt{torch.profiler}. We also count \texttt{nn.Module} calls with forward
hooks. Finally, we time batch-1 inference after compiling each model with \texttt{torch.compile} (default
Inductor backend), using the same repeat protocol. These sweep timings are separate from the batch-1,000 inference
times in \cref{tab:main}, which are measured on the trained models inside the training study; the two agree where
they overlap.

\subsection{Environment}
All main-study and sweep measurements run on a single machine (Apple M5 CPU, macOS, Python~3.9,
PyTorch~\torchVersion) with \texttt{torch.set\_num\_threads(1)} unless stated otherwise, so the timings
compare architectures rather than thread scheduling. No GPU is used: at this model scale, kernel-launch
overhead would dominate any GPU measurement~\cite{sze2017efficient}. \ifanonymous All code and raw results will be released publicly upon
publication.\else All code and raw results are available at
\url{https://github.com/aryanamol10/Research-Paper-Comparative-Analysis-of-Neural-Models}.\fi
